\begin{abstract}
Autoregressive vision-language models formulate medical question answering as language generation, even when the output is a choice from a predefined candidate set. We investigate direct multimodal decision modeling with Laya-Med, a System One-inspired non-autoregressive architecture that maps medical text and images to a distribution over candidate answers without generating answer tokens. Biomedical encoders, lightweight multimodal fusion, and a shared candidate-scoring head support batched decision inference. We organize the evaluation by backbone architecture and decision mechanism, comparing autoregressive generation, candidate-label likelihood scoring, and direct multimodal prediction across seven medical question-answering datasets. The study examines decision accuracy, probabilistic reliability, selective prediction, and observed inference cost. Matched pilot evaluations reveal task-dependent strengths rather than uniform superiority, while a separate full-cohort evaluation characterizes performance across heterogeneous medical tasks. A learned direct head on an autoregressive VLM is specified as a further control but awaits evaluation. The current evidence establishes a reproducible framework for studying direct medical decisions; matched architectural ablations and comprehensive memory comparisons remain future experiments.
\end{abstract}
